\section{Related Work}
\label{sec:related}

\paragraph{Network formation as discrete choice.}
Preferential attachment explains heavy-tailed degree distributions by new nodes linking to existing nodes in
proportion to their degree \citep{price1976general,barabasi1999emergence}; fitness models add a node-specific
attractiveness that lets late entrants overtake incumbents \citep{bianconi2001competition,wang2013quantifying}, and
empirical work measures the attachment kernel from the timing of new links
\citep{newman2001clustering,jeong2003measuring}. \citet{overgoor2019choosing} show that these mechanisms are
conditional-logit models in which each new edge is a choice among the existing nodes, so that degree, fitness,
homophily and triadic closure can be estimated jointly and compared by likelihood. With millions of alternatives
the choice set must be sampled; McFadden's correction for sampled alternatives \citep{mcfadden1978modelling} and its
importance-sampling extensions \citep{overgoor2020scaling} keep the estimates consistent. We build on this framework
and make two changes that matter at the scale of a model hub: candidate sets are drawn with probability
proportional to $(d+1)$, so that the competitors a developer actually weighs are almost always present, and the
mechanisms are compared with a Shapley decomposition of the likelihood gain instead of a sequence of nested models.

\paragraph{Popularity, herding and quality.}
When quality is hard to observe, adopters rationally follow earlier adopters
\citep{banerjee1992simple,bikhchandani1992theory}, and social influence makes success both more unequal and less
predictable from quality \citep{salganik2006experimental}. Cumulative advantage in science has the same structure
\citep{merton1968matthew}. Separating popularity from quality in observational data is the central difficulty of
this literature: popularity is itself a signal of quality. We address it in two ways: by bounding quality with
hub-specific fixed effects, the largest share any unobserved node attribute could take, and by measuring how much of
that bound external benchmark scores can account for.

\paragraph{Model ecosystems and model selection.}
Foundation models are adapted rather than trained from scratch \citep{bommasani2021opportunities}, which makes the
choice of a base model a consequential and recurring decision. Descriptive studies of the Hugging Face Hub document
its growth, its family trees and the concentration of downstream use on a few model families
\citep{osborne2024ai,laufer2025anatomy}. In machine learning, base-model choice is usually framed as a prediction
problem: transferability estimators such as LEEP and LogME rank pre-trained models by how well they are expected to
fine-tune on a target task \citep{nguyen2020leep,you2021logme}. Our question is complementary and behavioral: not
which model a developer \emph{should} choose, but which signals the choices developers actually make respond to,
and what kind of network those choices grow.
